\section{Applications to group von Neumann algebras}

We will now turn our attention to group von Neumann algebras. Let $G$ be a group. A sequ\-ence~$(g_n)_n$ of elements in $G$ is called asymptotically free if $(g_n)_n$ is free from the diagonal copy of $G$ in $\prod_n G$. Equivalently, for every $k \in \mathbb N$, $s_1,\dots,s_k \in G \setminus \{1\}$ and every $\varepsilon_1,\dots,\varepsilon_k \in \mathbb Z \setminus \{0\}$, there exists $n \in \mathbb N$, such that
\begin{gather*}
s_1 g_n^{\varepsilon_1} s_2 g_n^{\varepsilon_2} s_3 \cdots g_n^{\varepsilon_k} \neq 1.
\end{gather*}
Groups admitting an asymptotically free sequence have been studied for a long time.

Recall that a mixed identity for $G$ is a word $w \in \mathbb Z \ast G = \langle G, t\rangle$ in one variable $t$ with coefficients in $G$ such that $w$ evaluates to the identity if any element of $G$ is substituted for $t$. A~group is called mixed identity free (MIF) if it does not satisfy any non-trivial mixed identity. We~refer the reader to \cite{MR3556970} for further details around this notion, emphasizing the following statement.

\begin{Proposition}[{\cite[proof of Proposition 5.3 and Remark 5.1]{MR3556970}}]
Let $G$ be a group. The following are equivalent:
\begin{enumerate}\itemsep=0pt
\item[$1.$] The group $G$ contains an asymptotically free sequence.
\item[$2.$] There is no mixed identity $w \in \mathbb Z \ast G$ for $G$.
\item[$3.$] $G$ is mixed identity free, i.e., there exists no mixed identity in $\mathbb F_\infty \ast G$.
\end{enumerate}
\end{Proposition}

We denote the group von Neumann algebra of $G$ with its natural trace by $(LG,\tau)$. It is well-known that any non-trivial conjugacy class in a ${\rm MIF}$~group is infinite. In particular, the group von Neumann algebra $LG$ of a ${\rm MIF}$~group $G$ is a ${\rm II}_1$-factor. Note that the left-regular representation provides a natural inclusion $G \subset {\rm PU}(LG)$.
 We will be interested in the existence if maximal discrete subgroups of ${\rm PU}(LG)$ containing $G$.

\begin{Lemma}
Let $(g_n)_n$ be an asymptotically free sequence. There exists an ultrafilter $\omega$, such that $[(g_n)_n] \in U((LG)^{\omega})$ is freely independent from the diagonal copy of $LG$ in $(LG)^\omega$.
\end{Lemma}
This lemma is standard but not completely trivial, and we include the argument for the sake of completeness.
\begin{proof}
Let $(w_m)$ be an enumeration of all mixed identities, that is, elements in $\mathbb Z\ast G$. We~denote by $w(m_1,\dots,m_\ell)$ the iterated commutator
\begin{gather*}
w(m_1,\dots,m_\ell) = [w_{m_1},[w_{m_2},\dots, [w_{m_{\ell-1}},w_{m_\ell}]]].
\end{gather*}
Observe that if $w(m_1,\dots,m_\ell)(g) \neq 1$, it guarantees that $w_{m_i}(g)\neq 1$ for all $i=1,\dots,\ell$. Let $v_\ell\coloneqq w(1,2,\dots,\ell)$. Asymptotic freeness of $(g_n)$ guarantees the existence of natural numbers $n_\ell$ such that $v_\ell(g_{n_\ell})\neq 1$. Taking an ultrafilter containing this sequence finishes the proof.
\end{proof}

\begin{Lemma}
Let $G$ be a ${\rm MIF}$~group and $u \in {\rm U}(LG)$ with $u \not \in S^1 \cdot 1$ and $\operatorname{Re}(\tau(u))>3/4$. For any $\varepsilon>0$ there exists some $w \in \mathbb F_2$ and $g \in G$ such that $\ell(w(u,g))< \varepsilon$. Moreover, the image of the map from $\mathbb Z \ast G$ to ${\rm PU}(LG)$ that maps the generator $1 \in \mathbb Z$ to $u$ is not discrete.
\end{Lemma}
\begin{proof}
Note that from the assumptions $\bar \ell(u) \neq 0$ and $\ell(u)< \sqrt{2}$.
Let $(g_k)_k$ be an asymptotically free sequence and set $v= [(g_k)_k] \in (LG)^{\omega}$. By assumption $\ell(u) \leq 1/\sqrt{2} - \delta$ for some $\delta>0$. By~the results from above, we get
\begin{gather*}
\big(1/\sqrt{2}\big)^{n-1} \cdot \bar\ell(u)^n \leq \ell(w_n(u,v)) \leq\big(\sqrt{2}\big)^{n-1} \cdot \ell(u)^n.
\end{gather*}
Thus, for some $k \in \mathbb N$
\begin{gather*}
\big(1/\sqrt{2}\big)^{n} \cdot \bar\ell(u)^n \leq \ell(w_n(u,g_k)) \leq \big(\sqrt{2}\big)^{n} \cdot \ell(u)^n
\end{gather*}
and this finishes the proof.
\end{proof}

The following consequence is immediate.
\begin{Corollary}
Let $G$ be a ${\rm MIF}$~group and let $u \in {\rm PU}(LG)$ with $0 < \bar \ell(u) < 1/\sqrt{2}$. Then, the subgroup $\langle u, G \rangle \subset {\rm PU}(LG)$ is not discrete.
In particular, the set of discrete subgroups of~${\rm PU}(LG)$ containing $G$ is uniformly discrete.
\end{Corollary}
\begin{proof} By the previous lemma, no such subgroup contains a non-trivial element with $\bar \ell(u)<1/\sqrt{2}$. Indeed, we would apply the previous lemma to a suitable lift $u'$ with $\ell(u')< 1/\sqrt{2}$.\end{proof}

We are now able to state and prove our main result.
\begin{Theorem} %\label{main}
Let $G$ be a ${\rm MIF}$~group. Then, $G$ is contained in a maximal discrete subgroup in~${\rm PU}(LG)$.
\end{Theorem}
\begin{proof}
The discrete subgroups of ${\rm PU}(LG)$ containing $G$ are ordered by inclusion. Now, the result follows by Zorn's lemma since unions over chains remain discrete.
\end{proof}
